\documentclass{article}
\usepackage{/globals/anki}
\ankiprefix{crm-systeme}
\ankitopic{Customer-Relationship-Management}
\ankishow
\begin{document}
\section{Customer-Relationship-Management} 
Beim Customer-Relationship-Management (CRM) werden versucht Kunden zu behalten. Dabei werden deren Daten gespeichert und sie werden aktiv angesprochen.
\begin{description}
 \item[Kommunikatives CRM] Es wird Kundenkontakt gehalten.
 \item[Operatives CRM] Automatisierung vom Marketing, direkte, individuelle Ansprache. Daten werden gesammelt.
 \item[Analytisches CRM] Datamining von Kundenbedürfnissen.
\end{description}
\ankinote{aufgabe}{Aufgabe}{Kunden zu behalten, Kontakt}
\ankinote{arten}{Arten}{Kommunikativ, Operativ, Analytisch}
\ankinote{komm}{Kommunikativ}{Kundenkontakt}
\ankinote{ops}{Operativ}{Automatisierung Marketing, Ansprache, Datensammlung}
\ankinote{analytisch}{Analytisch}{Datamining}
\end{document}